% Lösungen Einheit 3 von 4 – Normalform und p-q-Formel (zusammenbau v0.1)
\einheitenkopf{Einheit 3 von 4 · Normalform und p-q-Formel}

\erg{16}{a) Normalform – Formel}
\erg{17}{a) durch $4$ – die Vorzahl von $x^2$}
%% TODO Lösung der Erklärzeile 18g) fehlt in der Bank
\erg{18}{a) $p = -1$, $q = -12$ \quad b) $x_{1,2} = -3 \pm \sqrt{9 - 8} = -3 \pm 1$; $x_1 = -2$, $x_2 = -4$ \quad c) $x_{1,2} = -5 \pm \sqrt{25 - 16} = -5 \pm 3$; $x_1 = -2$, $x_2 = -8$ \quad d) $x_{1,2} = -6 \pm \sqrt{36 - 20} = -6 \pm 4$; $x_1 = -2$, $x_2 = -10$ \quad e) $x_{1,2} = -7 \pm \sqrt{49 - 24} = -7 \pm 5$; $x_1 = -2$, $x_2 = -12$ \quad f) $x_{1,2} = -5 \pm \sqrt{25 - 21} = -5 \pm 2$; $x_1 = -3$, $x_2 = -7$ \quad g) \ldots \quad h) $x_{1,2} = 2 \pm \sqrt{4 + 12} = 2 \pm 4$; $x_1 = 6$, $x_2 = -2$ \quad i) $x^2 - 6x + 5 = 0$; $x_{1,2} = 3 \pm \sqrt{9 - 5} = 3 \pm 2$; $x_1 = 5$, $x_2 = 1$ \quad j) durch $5$ geteilt: $x^2 - 4x + 3 = 0$; $x_{1,2} = 2 \pm \sqrt{4 - 3} = 2 \pm 1$; $x_1 = 3$, $x_2 = 1$ \quad k) $x_{1,2} = 1{,}5 \pm \sqrt{2{,}25 - 2} = 1{,}5 \pm 0{,}5$; $x_1 = 2$, $x_2 = 1$ \quad l) $x_{1,2} = 5 \pm \sqrt{25 - 25} = 5$ – genau eine Lösung: $x = 5$ \quad m) Unter der Wurzel: $1 - 5 = -4$, negativ: keine Lösung \quad n) $D = 16 - 7 = 9$, positiv: zwei Lösungen \quad o) $x_{1,2} = 3 \pm \sqrt{9 - 4} = 3 \pm \sqrt{5}$; $x_1 \approx 5{,}24$, $x_2 \approx 0{,}76$ \quad p) $(-2)^2 - 3 \cdot (-2) - 10 = 4 + 6 - 10 = 0$ (wA): ja \quad q) $x^2 + 4x + 4 = 5x + 10$; $x^2 - x - 6 = 0$; $x_{1,2} = 0{,}5 \pm \sqrt{0{,}25 + 6} = 0{,}5 \pm 2{,}5$; $x_1 = 3$, $x_2 = -2$ \quad r) Nullprodukt: $x \cdot (x - 12) = 0$; $x_1 = 0$, $x_2 = 12$ \quad s) $2x^2 - 2x - 40 = 0 \mid : 2$; $x^2 - x - 20 = 0$; $x_{1,2} = 0{,}5 \pm \sqrt{0{,}25 + 20} = 0{,}5 \pm 4{,}5$; $x_1 = 5$, $x_2 = -4$}
\erg{19}{a) $x^2 = x + 6$; $x^2 - x - 6 = 0$; $x_{1,2} = 0{,}5 \pm \sqrt{0{,}25 + 6} = 0{,}5 \pm 2{,}5$; $x_1 = 3$, $x_2 = -2$}
\erg{20}{a) Nicht normiert – erst durch $3$ teilen. Richtig: $x^2 + 2x - 8 = 0$; $x_{1,2} = -1 \pm \sqrt{1 + 8} = -1 \pm 3$; $x_1 = 2$, $x_2 = -4$ \quad b) Die Formel gilt nur für die Normalform $x^2 + px + q = 0$ mit eins vor $x^2$; ohne Teilen sind p und q falsch abgelesen \quad c) Nullstellen ablesen: $x_1 = 3$, $x_2 = 1$ \quad d) $-0{,}5x^2 + 3x + 8 = 0 \mid : (-0{,}5)$; $x^2 - 6x - 16 = 0$; $x_{1,2} = 3 \pm \sqrt{9 + 16} = 3 \pm 5$; $x_1 = 8$, $x_2 = -2$; der Ball trifft nach $8$ m auf, $-2$ m fällt weg}
